\section{Metric Identity}

The real package records:

\[
  f_{\mathrm{cut}} =
  \frac{\text{tract adjacency edges crossing districts}}
       {\text{all tract adjacency edges}}.
\]

This metric ignores polygon perimeter, boundary length, projection, coastline,
and area. It is a graph-topology fragmentation measure. The benchmark itself
uses shared-boundary weights, so even its optimization objective is not
identical to the unweighted evaluation metric
~\citep{polsby1991third,reock1961note}.

\section{Observed Positions}

Rhode Island's benchmark is at the 1.35th percentile in Rust and 1.81st in
GerryChain. Iowa is at the 0.33rd and 0.88th percentiles, with ESS 257 and 498.
North Carolina is at 0 and 0.2 percent, with ESS 109 and 87. Under the
preregistered rule, no sub-1\% claim may be published unless ESS is at least
1,000. Thus only Rhode Island supports a stable bounded statement: the
benchmark lies near the low end of the sampled cut-fraction distribution.

\section{Implementation Sensitivity}

The Rust and GerryChain kernels use the same random-weight Kruskal tree method,
initial plan, population tolerance, and chain length, yet their cut
distributions differ. The KS effect is small in Rhode Island (0.038), moderate
in Iowa (0.162), and large in North Carolina (0.471). Proposal details,
pair reselection, random streams, and implementation behavior therefore matter
to distributional claims.

\section{Withdrawn Claims}

The package contains no polygon compactness trace. It cannot support:

\begin{itemize}
\item the former 61st--72nd Polsby--Popper percentile range;
\item the claim that no benchmark exceeds the 95th compactness percentile;
\item a compactness--partisan correlation; or
\item a legal ``optimal compactness position.''
\end{itemize}

Those questions require archived geometries or per-plan polygon metrics under a
versioned projection and metric implementation.

\section{Conclusion}

Metric identity is part of evidence identity. A low tract cut fraction is a
real, reproducible result for this package, but calling it ``Polsby--Popper
compactness'' would be a category error.
